\subsection{The Regular Polygon}
%

\begin{definition}
	In Fig. \ref{ch5_polygon_def}, 6 congruent triangles are arranged in a circular fashion.  Such a figure is known as a regular hexagon.  In general, $n$ number of traingles can be arranged to form a regular polygon.
\end{definition}
\begin{figure}[!h]
	\begin{center}
		
		\includegraphics[width=\columnwidth]{./figs/ch5_polygon_def}
		\vspace*{-10cm}
	\end{center}
	\caption{Polygon Definition}
	\label{ch5_polygon_def}	
\end{figure}
%
\begin{definition}
The angle formed by each of the congruent triangles at the centre of a regular polygon of $n$ sides is $\frac{360^{\degree}}{n}$.
\end{definition}
%
\begin{problem}
Show that the area of a regular polygon is given by 
%
\begin{equation}
\frac{n}{2}r^{2}\sin\frac{360^{\degree}}{n}
\end{equation}
%
\end{problem}
\begin{figure}[!h]
	\begin{center}
		
		\includegraphics[width=\columnwidth]{./figs/ch5_polygon_area}
		\vspace*{-10cm}
	\end{center}
	\caption{Polygon Area}
	\label{ch5_polygon_area}	
\end{figure}
%

\proof The triangle that forms the polygon of $n$ sides is given in Fig. \ref{ch5_polygon_area}.  Thus,
%
\begin{equation}
\begin{split}
ar\brak{polygon} &= n ar\brak{\Delta ABC} \\
&= \frac{n}{2}r^{2}\sin\frac{360^{\degree}}{n}
\end{split}
\end{equation}
%
\begin{problem}
	Using Fig. \ref{ch5_circle_squeeze}, show that
%
\begin{equation}
\label{ch5_circle_squeeze_eq}
\frac{n}{2}r^{2}\sin\frac{360^{\degree}}{n} < \text{ area of circle } < nr^{2}\tan\frac{180^{\degree}}{n}
\end{equation}
%
The portion of the circle visible in Fig. \ref{ch5_circle_squeeze} is defined to be a sector of the circle.
\end{problem}
\begin{figure}[!h]
	\begin{center}
		
		\includegraphics[width=\columnwidth]{./figs/ch5_circle_squeeze}
		\vspace*{-10cm}
	\end{center}
	\caption{Circle Area in between Area of Two Polygons}
	\label{ch5_circle_squeeze}	
\end{figure}
%

\proof Note that the circle is squeezed between the inner and outer regular polygons.  As we can see from Fig. \ref{ch5_circle_squeeze}, the area of the circle should be in between the areas of the inner and outer polygons.  Since
%
\begin{align}
ar \brak{\Delta OAB} &= \frac{1}{2}r^{2}\sin\frac{360^{\degree}}{n} \\
ar \brak{\Delta OPQ} &= 2 \times \frac{1}{2} \times r \tan\frac{360/n}{2} \times r \\
&= r^{2}\tan\frac{180^{\degree}}{n},
\end{align}
%
we obtain \eqref{ch5_circle_squeeze_eq}.
%
%
\begin{problem}
	Using Fig. \ref{ch5_sin_theta}, show that 
	%
\begin{equation}
\label{ch5_sin_theta_eq}
\sin  \theta_1 = \sin \brak{\theta_1 + \theta_2}\cos \theta_2 - \cos\brak{\theta_1+\theta_2}\sin\theta_2
\end{equation}	
	%
\end{problem}
\begin{figure}[!h]
	\begin{center}
		
		\includegraphics[width=\columnwidth]{./figs/ch5_sin_theta}
		\vspace*{-10cm}
	\end{center}
	\caption{$\sin2\theta = 2\sin\theta\cos\theta$}
	\label{ch5_sin_theta}	
\end{figure}
%

\proof The following equations can be obtained from the figure using the forumula for the area of a triangle
%
\begin{align}
ar \brak{\Delta ABC} &= \frac{1}{2}ac \sin\brak{\theta_1 + \theta_2} \\
&= ar \brak{\Delta BDC} + ar \brak{\Delta ADB} \\
&= \frac{1}{2}cl \sin{\theta_1} + \frac{1}{2}al \sin{\theta_2} \\ 
&= \frac{1}{2}ac \sin{\theta_1} \sec \theta_2 + \frac{1}{2}a^2 \tan{\theta_2} 
\end{align}
$\brak{\because
	l = a \sec \theta_2}$.  From the above,
\begin{align}
\Rightarrow \sin\brak{\theta_1 + \theta_2} &=  \sin{\theta_1} \sec \theta_2 + \frac{a}{c} \tan{\theta_2} \\
\Rightarrow \sin\brak{\theta_1 + \theta_2} &=  \sin{\theta_1} \sec \theta_2 + \cos\brak{\theta_1 + \theta_2} \tan{\theta_2} 
\end{align}
Multiplying both sides by $\cos \theta_2$,
\begin{align}
\Rightarrow \sin\brak{\theta_1 + \theta_2}\cos{\theta_2} &=  \sin{\theta_1}  + \cos\brak{\theta_1 + \theta_2} \sin\theta_2  
\end{align}
%
resulting in
\begin{equation}
\Rightarrow \sin \theta_1 = \sin\brak{\theta_1 + \theta_2}\cos{\theta_2} - \cos\brak{\theta_1 + \theta_2} \sin\theta_2 
\end{equation}
\begin{problem}
	Prove the following identities 
	%
	\begin{enumerate}
\item 
\begin{equation}
		\label{ch5_sin_diff}
\sin\brak{\alpha - \beta} = \sin \alpha \cos \beta - \cos \alpha \sin \beta.
\end{equation}
\item 
\begin{equation}
\cos\brak{\alpha + \beta} = \cos \alpha \cos \beta - \sin \alpha \sin \beta.
		\label{ch5_cos_diff}
\end{equation}

	\end{enumerate}
	%
\end{problem}
\proof In \eqref{ch5_sin_theta_eq}, let
%
\begin{equation}
\begin{split}
\theta_1 + \theta_2 &= \alpha \\
\theta_2 &=  \beta
\end{split}
\end{equation}
%
This gives \eqref{ch5_sin_diff}.  In \eqref{ch5_sin_diff}, replace $\alpha$ by 
%
$90^{\degree} - \alpha$.  This results in
%
\begin{multline}
\sin\brak{90^{\degree} - \alpha - \beta}
\\
=
\sin \brak{90^{\degree} -\alpha} \cos \beta - \cos \brak{90^{\degree} -\alpha} \sin \beta \\
\Rightarrow \cos\brak{\alpha + \beta} = \cos \alpha \cos \beta - \sin \alpha \sin \beta
\end{multline}
% 
\begin{problem}
	Using \eqref{ch5_sin_theta_eq} and \eqref{ch5_cos_diff}, show that
\begin{align}
\label{ch5_sin_sum}
\sin\brak{\theta_1 + \theta_2} &= \sin\theta_1  \cos\theta_2 + \cos\theta_1\sin\theta_2
\\
\cos\brak{\theta_1 - \theta_2} &= \cos\theta_1  \cos\theta_2  \sin\theta_1\sin\theta_2
\label{ch5_cos_sum}
\end{align}
\end{problem}
%
\proof From \eqref{ch5_sin_theta_eq},
%
\begin{align}
 \sin \brak{\theta_1 + \theta_2}\cos \theta_2 =\sin  \theta_1 +\cos\brak{\theta_1+\theta_2}\sin\theta_2 
\end{align}
%
Using \eqref{ch5_cos_diff} in the above,
%
\begin{multline}
\sin \brak{\theta_1 + \theta_2}\cos \theta_2 
=\sin  \theta_1 +\lbrak{\cos \theta_1\cos\theta_2 }
\\	
\rbrak{	- \sin \theta_1\sin\theta_2}\sin\theta_2 
\end{multline}
%
which can be expressed as
%
\begin{multline}
\sin \brak{\theta_1 + \theta_2}\cos \theta_2 
=\sin  \theta_1 +\cos \theta_1\cos\theta_2 \sin\theta_2 
\\	
	- \sin \theta_1\sin^2\theta_2
\end{multline}
%
Since
%
\begin{equation}
\sin^2\theta_2 = 1- \cos^2\theta_2, 
\end{equation}
%
we obtain
%
\begin{multline}
\sin \brak{\theta_1 + \theta_2}\cos \theta_2 
=\cos \theta_1\cos\theta_2 \sin\theta_2 
\\	
+ \sin \theta_1\cos^2\theta_2
\end{multline}
%
resulting in
%
\begin{equation}
\sin \brak{\theta_1 + \theta_2}
=\cos \theta_1 \sin\theta_2 
+ \sin \theta_1\cos\theta_2
\end{equation}
%
after factoring out $\cos \theta_2$.  Using a similar approach, \eqref{ch5_cos_sum} can also be proved.
%
\begin{problem}
	Show that
	%
	\begin{equation}
	\label{eq:sin2theta}
	\sin 2\theta = 2 \sin\theta \cos\theta
	\end{equation}
	%
\end{problem}
\begin{problem}
Show that
	%
	\begin{equation}
	\label{ch5_circle_squeeze_simple}
\cos^2\frac{180^{\degree}}{n} < \frac{\text{ area of circle }}{nr^{2}\tan\frac{180^{\degree}}{n}} < 1	\end{equation}
	%
\end{problem}
\proof From \eqref{ch5_circle_squeeze_eq} and \eqref{eq:sin2theta},
	%
	\begin{multline}
	\frac{n}{2}r^{2}\sin\frac{360^{\degree}}{n} < \text{ area of circle } 
	\\
	< nr^{2}\tan\frac{180^{\degree}}{n} 
	\\
\Rightarrow 	
	{n}r^{2}\sin\frac{180^{\degree}}{n}\cos\frac{180^{\degree}}{n} < \text{ area of circle } 
	\\
	< nr^{2}\tan\frac{180^{\degree}}{n} 
	\end{multline}
	%
\begin{problem}
	Show that if
\begin{equation}
\label{ch5_sin_increasing}
\theta_1 < \theta_2, \sin \theta_1 < \sin \theta_2.
\end{equation}	
\end{problem}
\proof Trivial using the definition and choosing angles $\theta_1$ and $\theta_2$ appropriately in a right angled triangle.
%
	%
	\begin{problem}
		Show that if
		\begin{equation}
		\label{ch5_sin_increasing}
		\theta_1 < \theta_2, \cos \theta_1 > \cos \theta_2.
		\end{equation}	
	\end{problem}
	%
\begin{problem}
	Show that 
	\begin{equation}
	\label{ch5_sin_zero}
	\sin 0^{\degree} = 0
	\end{equation}
\end{problem}
\proof Follows from \eqref{ch5_sin_increasing}.
%
\begin{problem}
	Show that 
	\begin{equation}
	\label{ch5_sin_zero}
	\cos 0^{\degree} = 1
	\end{equation}
	\end{problem}
%
%
\begin{problem}
	Show that for large values of $n$
	%
	\begin{equation}
	%
\cos^2\frac{180^{\degree}}{n} = 1
%
	\end{equation}	
	% 
\end{problem}
%
\proof  Follows from previous problem.
%
\begin{definition}
	The previous result can be expressed as
%
\begin{equation}
\lim_{n \rightarrow \infty}\cos^2\frac{180^{\degree}}{n} = 1
\end{equation}
%	
\end{definition}
\begin{problem}
	Show that 
	%
\begin{equation}
\text{ area of circle } = r^2\lim_{n \rightarrow \infty}
{n\tan\frac{180^{\degree}}{n}} 
	%
\end{equation}	
	% 
\end{problem}
%
\begin{definition}
	\begin{equation}
	\pi = \lim_{n \rightarrow \infty}
	{n\tan\frac{180^{\degree}}{n}}
	\end{equation}
\end{definition}
Thus, the area of a circle is $\pi r^2$.
\begin{definition}
	The radian is a unit of angle defined by
\begin{equation}
	1 \text{ radian} = \frac{360^{\degree}}{2\pi}
\end{equation}
\end{definition}
%
\begin{problem}
	Show that the circumference of a circle is $2 \pi r$.
\end{problem}
\begin{problem}
	Show that the area of a sector with angle $\theta$ in radians is $\frac{1}{2}r^2\theta$.
\end{problem}
\begin{problem}
	Show that
	\begin{equation}
	\lim_{\theta \rightarrow 0} \frac{\sin\theta}{\theta} = 1
	\end{equation}
\end{problem}
